\documentclass[12pt, a4paper]{article}

% --- PACOTES ---
\usepackage[utf8]{inputenc}
\usepackage[T1]{fontenc}
\usepackage[brazil]{babel}
\usepackage{amsmath, amssymb, amsfonts}
\usepackage{geometry}
\usepackage{graphicx}
\usepackage{indentfirst}
\usepackage{setspace}
\usepackage{titlesec}
\usepackage{enumitem}
\usepackage{float} % Para posicionar tabelas e figuras

% --- CONFIGURAÇÃO DA PÁGINA ---
\geometry{left=3cm, right=2cm, top=3cm, bottom=2cm}
\setlength{\parskip}{6pt}
\onehalfspacing

% --- DADOS DO TRABALHO ---
\title{\textbf{ANÁLISE APROFUNDADA DE MÉTODOS DE OTIMIZAÇÃO LINEAR:}\\ \large DO SIMPLEX PADRÃO À DECOMPOSIÇÃO EM LARGA ESCALA}
\author{ÍGOR J. RODRIGUES, LEONARDO S. CORADELI, \\ LUCAS V. C. IKEDA, MARCO V. M. FARIA}
\date{2025}

\begin{document}

% --- CAPA ---
\begin{titlepage}
    \centering
    \textbf{\Large UNIVERSIDADE ESTADUAL PAULISTA "JÚLIO DE MESQUITA FILHO"}\\
    \textbf{\large FACULDADE DE CIÊNCIA E TECNOLOGIA (FCT)}
    
    \vspace{6cm}
    
    \textbf{\Large ANÁLISE APROFUNDADA DE MÉTODOS DE OTIMIZAÇÃO LINEAR:}\\
    \vspace{0.5cm}
    \large DO SIMPLEX PADRÃO À DECOMPOSIÇÃO EM LARGA ESCALA
    
    \vspace{5cm}
    
    \begin{flushright}
        \begin{minipage}{0.5\textwidth}
            Relatório técnico apresentado à disciplina de Pesquisa Operacional, como requisito parcial para avaliação.
        \end{minipage}
    \end{flushright}
    
    \vfill
    \textbf{Presidente Prudente}\\
    \textbf{2025}
\end{titlepage}

% --- SUMÁRIO ---
\tableofcontents
\newpage

% --- INÍCIO DO CONTEÚDO ---

\section{Introdução e Contextualização Histórica}
A Programação Linear (PL) representa um dos maiores triunfos da matemática aplicada do século XX. Desde a formulação do problema da dieta por George Stigler até o desenvolvimento do algoritmo Simplex por George Dantzig em 1947, a capacidade de otimizar uma função objetivo linear sujeita a restrições lineares revolucionou a economia, a logística e a engenharia.

No entanto, a complexidade dos problemas do mundo real cresceu exponencialmente nas últimas décadas. O que antes eram problemas manuais de planejamento de produção com dezenas de variáveis, hoje são problemas de otimização de redes energéticas ou escalas de tripulação aérea envolvendo milhões de variáveis e centenas de milhares de restrições.

Neste cenário, o Simplex Padrão ("Tableau Simplex"), embora didaticamente excelente, mostra-se inadequado para computação em larga escala. Este relatório explora as razões matemáticas para essa inadequação e apresenta as soluções robustas desenvolvidas pela comunidade científica: o Simplex Revisado e a Decomposição de Dantzig-Wolfe.

\section{Fundamentos Matemáticos e o Simplex Padrão}
Para compreender as vantagens do método revisado, é imperativo primeiro formalizar a álgebra do método padrão e identificar onde ocorrem seus gargalos computacionais.

\subsection{A Forma Padrão e a Geometria}
Considere o problema de Programação Linear na sua forma padrão:
\begin{equation}
    \begin{aligned}
        \text{Maximizar } \quad & z = c^T x \\
        \text{Sujeito a } \quad & Ax = b \\
                                & x \ge 0
    \end{aligned}
\end{equation}
Onde:
\begin{itemize}
    \item $A$ é uma matriz de restrições de dimensão $m \times n$ (com $n > m$ e posto completo $m$).
    \item $c \in \mathbb{R}^n$ é o vetor de custos.
    \item $b \in \mathbb{R}^m$ é o vetor de recursos (lado direito).
    \item $x \in \mathbb{R}^n$ é o vetor de variáveis de decisão.
\end{itemize}

O conjunto de soluções viáveis $P = \{x | Ax=b, x \ge 0\}$ forma um poliedro convexo em $\mathbb{R}^n$. O Teorema Fundamental da Programação Linear garante que, se existe uma solução ótima finita, ela reside em um dos vértices (pontos extremos) deste poliedro.

\subsection{Partição Básica e Soluções Viáveis}
A estratégia algébrica do Simplex baseia-se na partição da matriz $A$. Escolhemos $m$ colunas linearmente independentes para formar uma base $B$. As demais $n-m$ colunas formam a matriz não-básica $N$. Analogamente, particionamos as variáveis $x$ em $(x_B, x_N)$ e os custos $c$ em $(c_B, c_N)$.

A equação de restrições $Ax=b$ pode ser reescrita como:
\begin{equation}
    Bx_B + Nx_N = b
\end{equation}
Pré-multiplicando por $B^{-1}$, isolamos as variáveis básicas:
\begin{equation}
    x_B = B^{-1}b - B^{-1}Nx_N
\end{equation}
Uma Solução Básica Viável (SBV) ocorre quando fixamos $x_N = 0$. Logo, $x_B = B^{-1}b$. Se $x_B \ge 0$, estamos em um vértice do poliedro viável.

A função objetivo pode ser reescrita substituindo $x_B$:
\begin{align}
    z &= c_B^T x_B + c_N^T x_N \\
    z &= c_B^T (B^{-1}b - B^{-1}Nx_N) + c_N^T x_N \\
    z &= c_B^T B^{-1}b + (c_N^T - c_B^T B^{-1}N)x_N
\end{align}
O termo entre parênteses é o vetor de \textbf{Custos Reduzidos}, denotado por $\hat{c}_N^T$.

\subsection{O Mecanismo do Simplex via Tableau (Padrão)}
No método Simplex Padrão, mantemos e atualizamos iterativamente uma tabela (tableau) que representa explicitamente o sistema de equações lineares transformado pela base atual. Um tableau típico contém explicitamente a matriz $B^{-1}A$.

A cada iteração, realizamos operações de linha (pivoteamento de Gauss-Jordan) para transformar uma coluna não-básica em um vetor unitário da base canônica.

\subsubsection{Análise de Complexidade do Simplex Padrão}
Aqui reside a ineficiência crítica. Suponha um problema onde $m=1.000$ (restrições) e $n=10.000$ (variáveis).

\begin{enumerate}
    \item \textbf{Atualização Desnecessária:} Ao realizar o pivoteamento, o Simplex Padrão atualiza todas as $n$ colunas da matriz. Ou seja, calculamos $B^{-1}a_j$ para todas as variáveis não-básicas $j$, mesmo aquelas que têm custo reduzido muito negativo e provavelmente nunca entrarão na base.
    
    \item \textbf{Custo Computacional:} Uma operação de pivô requer aproximadamente $m \times n$ multiplicações e adições. Em problemas de grande porte, esse custo torna-se proibitivo.
    
    \item \textbf{Preenchimento (Fill-in):} Em problemas reais, a matriz $A$ é frequentemente esparsa (contém muitos zeros). Ao multiplicar $B^{-1}A$ explicitamente e fazer operações de linha, os zeros originais tendem a se tornar números não-nulos (fill-in). Isso aumenta drasticamente o uso de memória e o esforço computacional.
\end{enumerate}

\section{O Método Simplex Revisado: Uma Abordagem Eficiente}
O Método Simplex Revisado altera fundamentalmente a forma como os cálculos são realizados, mantendo a mesma trajetória geométrica (sequência de vértices) do método padrão.

\subsection{A Filosofia da Computação Sob Demanda}
A premissa do método revisado é: não precisamos ver o tableau inteiro para tomar a decisão de qual vértice visitar a seguir. Precisamos apenas de:
\begin{itemize}
    \item Os custos reduzidos ($\hat{c}$) para decidir quem entra na base.
    \item A coluna atualizada da variável de entrada ($\alpha$) para decidir quem sai (razão mínima).
\end{itemize}
Para isso, em vez de transformar a matriz $A$ inteira (uma abordagem de "Força Bruta"), trabalhamos sempre com a matriz $A$ original e mantemos apenas a inversa da base $B^{-1}$ atualizada (uma abordagem "Cirúrgica").

\subsection{O Algoritmo Matemático Revisado Detalhado}
Vamos detalhar o passo a passo matemático, enfatizando as diferenças de cálculo em relação ao padrão.

\subsubsection{Passo 1: Cálculo dos Multiplicadores Simplex (Vetor Dual)}
Em vez de calcular o custo reduzido diretamente, primeiro calculamos o vetor $\pi$ (multiplicadores simplex ou variáveis duais):
\begin{equation}
    \pi^T = c_B^T B^{-1}
\end{equation}
Isso equivale a resolver o sistema linear $B^T \pi = c_B$.

\subsubsection{Passo 2: Pricing (Seleção da Variável de Entrada)}
Para encontrar a variável de entrada, calculamos:
\begin{equation}
    \hat{c}_j = c_j - \pi^T a_j \quad \forall j \in \text{Não-Básicas}
\end{equation}
Note a eficiência: este cálculo é um produto interno entre vetores. Não precisamos calcular a coluna inteira atualizada. Além disso, podemos usar estratégias de \textit{Pricing Parcial}, calculando $\hat{c}_j$ apenas para um subconjunto de variáveis até encontrar uma candidata positiva, economizando CPU.

Se $\hat{c}_j \le 0$ para todo $j$, PARE. A solução atual é ótima. Caso contrário, selecione $k$ tal que $\hat{c}_k > 0$.

\subsubsection{Passo 3: Cálculo da Coluna Atualizada}
Somente agora, e apenas para a variável vencedora $x_k$, calculamos sua representação na base atual:
\begin{equation}
    \alpha = B^{-1} a_k
\end{equation}
Isso equivale a resolver o sistema linear $B\alpha = a_k$.

\subsubsection{Passo 4: Teste da Razão (Variável de Saída)}
Idêntico ao método padrão. Calculamos o vetor de solução atual $\bar{b} = B^{-1}b$ (normalmente já disponível da iteração anterior) e fazemos:
\begin{equation}
    \theta = \min \left\{ \frac{\bar{b}_i}{\alpha_i} : \alpha_i > 0 \right\}
\end{equation}
Seja $x_{B_r}$ a variável básica associada ao índice $r$ que minimiza a razão. Ela deixará a base.

\subsubsection{Passo 5: Atualização da Inversa da Base}
Este é o coração do método revisado. Temos uma nova base $B_{nova}$ que difere de $B_{velha}$ por apenas uma coluna. Não precisamos inverter a matriz do zero (custo $O(m^3)$). Usamos matrizes elementares. A nova inversa é dada por:
\begin{equation}
    B_{nova}^{-1} = E \cdot B_{velha}^{-1}
\end{equation}
Onde $E$ é uma matriz elementar (Matriz Eta) que transforma a coluna $r$-ésima da identidade na coluna pivô.

\subsection{Vantagens Computacionais Críticas}

\subsubsection{Preservação da Esparsidade}
Como a matriz $A$ nunca é alterada, seus zeros permanecem zeros. Isso economiza imensa quantidade de memória. Apenas a matriz $B^{-1}$ (ou sua representação fatorada) precisa ser armazenada.

\subsubsection{Estabilidade Numérica}
No Simplex Padrão, erros de arredondamento se propagam cumulativamente por todas as $mn$ entradas do tableau a cada iteração. Após milhares de iterações, o tableau pode se tornar "lixo numérico".

No Revisado, os erros se acumulam apenas na inversa da base. Periodicamente, podemos descartar a inversa acumulada e recalcular $B^{-1}$ do zero usando decomposição LU de alta precisão sobre a matriz original. Isso "reseta" o erro numérico, uma funcionalidade impossível no método padrão.

\subsubsection{Comparativo de Operações}
\begin{itemize}
    \item \textbf{Padrão:} Custo de $O(m \cdot n)$ por iteração.
    \item \textbf{Revisado:} Custo de $O(m^2)$ por iteração (assumindo cálculo denso).
\end{itemize}
Conclusão: Quando $n \gg m$ (muito mais variáveis que restrições), o Revisado é ordens de grandeza mais rápido.

\section{Exemplo Prático de Aplicação do Simplex Revisado}
Para consolidar a teoria apresentada, aplicaremos o algoritmo revisado passo a passo em um problema numérico.

\textbf{Problema:} Maximizar $Z = 4x_1 + 3x_2$
$$ \text{Sujeito a: } \begin{cases} 2x_1 + x_2 \le 10 \\ x_1 + x_2 \le 8 \\ x_1, x_2 \ge 0 \end{cases} $$

\textbf{Preparação (Forma Padrão):}
Introduzimos as variáveis de folga $x_3$ e $x_4$.
$$ \text{Max } Z = 4x_1 + 3x_2 + 0x_3 + 0x_4 $$
$$ \begin{cases} 2x_1 + 1x_2 + 1x_3 + 0x_4 = 10 \\ 1x_1 + 1x_2 + 0x_3 + 1x_4 = 8 \end{cases} $$

A matriz $A$, o vetor $b$ e o vetor $c$ são:
$$ A = \begin{bmatrix} 2 & 1 & 1 & 0 \\ 1 & 1 & 0 & 1 \end{bmatrix}, \quad b = \begin{bmatrix} 10 \\ 8 \end{bmatrix}, \quad c = \begin{bmatrix} 4 & 3 & 0 & 0 \end{bmatrix}^T $$

\textbf{Iteração 1:}
\begin{enumerate}
    \item \textbf{Base Inicial:} Escolhemos as folgas. $B = I = \begin{bmatrix} 1 & 0 \\ 0 & 1 \end{bmatrix}$, $B^{-1} = I$. Variáveis básicas $x_B = \{x_3, x_4\}$. Custos básicos $c_B = [0, 0]$.
    
    \item \textbf{Passo 1 (Multiplicadores $\pi$):}
    $$ \pi^T = c_B^T B^{-1} = [0, 0] \cdot I = [0, 0] $$
    
    \item \textbf{Passo 2 (Pricing):} Calculamos $\hat{c}_j = c_j - \pi^T a_j$.
    $$ \hat{c}_1 = 4 - [0,0]\cdot \binom{2}{1} = 4 $$
    $$ \hat{c}_2 = 3 - [0,0]\cdot \binom{1}{1} = 3 $$
    O maior custo reduzido positivo é $\hat{c}_1 = 4$. Logo, \textbf{$x_1$ entra na base}.
    
    \item \textbf{Passo 3 (Coluna Atualizada):}
    $$ \alpha = B^{-1} a_1 = I \cdot \binom{2}{1} = \binom{2}{1} $$
    
    \item \textbf{Passo 4 (Teste da Razão):}
    Solução atual $\bar{b} = B^{-1}b = [10, 8]^T$.
    $$ \text{Razão 1 ($x_3$): } 10 / 2 = 5 $$
    $$ \text{Razão 2 ($x_4$): } 8 / 1 = 8 $$
    A menor razão é 5. Logo, a variável básica da primeira linha ($x_3$) \textbf{sai da base}.
    
    \item \textbf{Resultado da Iteração 1:} Nova base será $\{x_1, x_4\}$.
\end{enumerate}

Este exemplo demonstra claramente que não foi necessário calcular a coluna atualizada de $x_2$ ou atualizar toda a matriz de restrições, economizando operações.

\section{Aplicações Práticas: Quando usar o Revisado?}
O Simplex Revisado é o padrão da indústria em solvers comerciais como CPLEX, Gurobi e XPRESS. Ele brilha em cenários específicos:

\subsection{Problemas de Escala de Tripulação (Crew Scheduling)}
Companhias aéreas precisam alocar milhares de tripulantes para milhares de voos.
\begin{itemize}
    \item \textbf{Estrutura:} Milhões de variáveis (possíveis combinações de rotas/escalas) e "apenas" alguns milhares de restrições (cobrir todos os voos, leis trabalhistas).
    \item \textbf{Por que o Revisado?} Com $n \approx 1.000.000$ e $m \approx 5.000$, manter um tableau de $5 \cdot 10^9$ entradas é impossível. O Revisado calcula custos reduzidos sob demanda, permitindo a resolução eficiente.
\end{itemize}

\subsection{Problemas Dinâmicos e Estocásticos}
Em problemas de múltiplos estágios onde novas variáveis são geradas dinamicamente (Geração de Colunas), o método revisado é obrigatório, pois a matriz $A$ completa nem sequer existe no início do processo.

\section{Princípio de Decomposição de Dantzig-Wolfe}
A evolução natural para resolver problemas que excedem até mesmo a capacidade do Simplex Revisado é a decomposição. O método de Dantzig-Wolfe aborda problemas com estrutura \textbf{Bloco-Angular}.

\subsection{Estrutura Bloco-Angular}
Imagine uma corporação com duas fábricas. Cada fábrica tem suas restrições internas (máquinas, turnos), mas ambas competem por um orçamento corporativo global. A matriz de restrições tem a forma:
\begin{equation}
    A = \begin{bmatrix}
    L_1 & L_2 & \dots & L_K \\
    A_1 & 0 & \dots & 0 \\
    0 & A_2 & \dots & 0 \\
    \vdots & \vdots & \ddots & \vdots \\
    0 & 0 & \dots & A_K
    \end{bmatrix}
\end{equation}
Onde $L_k$ são as restrições de ligação (globais) e $A_k$ são os blocos independentes.

\subsection{Reformulação do Problema Mestre e Convexidade}
A ideia central baseia-se na **Teoria da Convexidade** e no **Teorema da Representação de Poliedros**. Qualquer solução viável para o bloco $k$ ($x_k$) pode ser expressa como uma combinação convexa dos vértices do poliedro deste bloco.

Seja $V_k = \{v_{k1}, v_{k2}, \dots, v_{kp}\}$ o conjunto de vértices do poliedro definido por $A_k x_k = b_k$. Então:
\begin{equation}
    x_k = \sum_{j} \lambda_{kj} v_{kj}, \quad \sum \lambda_{kj} = 1, \quad \lambda \ge 0
\end{equation}
Substituímos $x_k$ no problema original pelas variáveis $\lambda$. O problema original nas variáveis $x$ (que tem muitas restrições) transforma-se em um **Problema Mestre** nas variáveis $\lambda$ (que tem poucas restrições globais, mas muitas variáveis — uma para cada vértice de cada subproblema).

\subsection{O Ciclo de Geração de Colunas}
Não podemos listar todos os vértices $v_{kj}$ a priori, pois são muitos. Usamos o Simplex Revisado no Problema Mestre e geramos as colunas (os vértices) sob demanda.

\begin{enumerate}
    \item \textbf{Mestre Restrito:} Resolvemos o problema mestre com apenas alguns vértices conhecidos. Obtemos os preços sombra duais ($\pi$) das restrições de ligação.
    
    \item \textbf{Subproblemas:} Enviamos $\pi$ para as divisões. Cada divisão $k$ resolve um subproblema local tentando minimizar seu custo reduzido ajustado pelo preço sombra:
    $$ \min (c_k - \pi^T L_k)x_k \quad \text{s.a.} \quad A_k x_k = b_k $$
    
    \item \textbf{Integração:} A solução ótima do subproblema é um vértice $v_{kj}^*$. Se ele tiver custo reduzido negativo no mestre, adicionamos essa nova coluna ao Mestre e reotimizamos.
\end{enumerate}

Este processo imita um mercado econômico: o centro (Mestre) define preços para recursos escassos, e as filiais (Subproblemas) propõem planos de produção baseados nesses preços.

\section{Conclusão}
A transição do Simplex Padrão para o Revisado não é apenas uma mudança de implementação, mas uma mudança paradigmática de "processamento em lote" para "processamento sob demanda". O Simplex Padrão é pedagogicamente inestimável para entender a geometria dos vértices, mas computacionalmente inviável para o mundo real devido à densidade e complexidade $O(mn)$.

O Simplex Revisado, ao explorar a álgebra matricial esparsa e focar nas operações essenciais ($B^{-1}$), viabiliza a resolução de problemas industriais. Para além disso, o Princípio de Decomposição de Dantzig-Wolfe eleva essa capacidade, permitindo que sistemas massivos e descentralizados sejam otimizados de forma coordenada, pavimentando o caminho para a moderna otimização em larga escala.

\section{Referências Bibliográficas}
\begin{enumerate}
    \item ARENALES, M.; ARMENTANO, V.; MORABITO, R.; YANASSE, H. \textit{Pesquisa Operacional: Para Cursos de Engenharia}. Rio de Janeiro: Elsevier, 2007.
    \item BAZARAA, M. S.; JARVIS, J. J.; SHERALI, H. D. \textit{Linear Programming and Network Flows}. 4. ed. New Jersey: John Wiley \& Sons, 2010.
    \item BERTSEKAS, D. P. \textit{Linear Network Optimization: Algorithms and Codes}. MIT Press, 1991.
    \item CHVÁTAL, V. \textit{Linear Programming}. New York: W. H. Freeman, 1983.
    \item DANTZIG, G. B. \textit{Linear Programming and Extensions}. Princeton University Press, 1963.
    \item LASDON, L. S. \textit{Optimization Theory for Large Systems}. New York: Macmillan, 1970.
    \item VANDERBEI, R. J. \textit{Linear Programming: Foundations and Extensions}. 4. ed. Springer, 2014.
\end{enumerate}

\end{document}